% ---------------------------------------------------------------------
%  CHƯƠNG 2 — CƠ SỞ LÝ THUYẾT VÀ CÔNG NGHỆ
% ---------------------------------------------------------------------
\chapter{Cơ sở lý thuyết và công nghệ}
\label{ch:coso}

\begin{muctieu}
Chương này trình bày các lý thuyết nền tảng về hệ chuyên gia y tế, hệ thống
phân tầng nguy cơ lâm sàng, mô hình ngôn ngữ lớn và kỹ thuật gọi hàm có cấu
trúc, đồng thời giới thiệu hệ sinh thái công nghệ được sử dụng để xây dựng giải pháp.
\end{muctieu}

\section{Hệ chuyên gia và suy diễn dựa trên luật}

Hệ chuyên gia \en{(Expert System)} là một trong những nhánh phát triển thành
công sớm nhất của trí tuệ nhân tạo mang tính biểu tượng \en{(Symbolic AI)}
\cite{RussellNorvig2021, HoangKiem}. Trong y khoa, các hệ thống như MYCIN \cite{Buchanan1984}
đã chứng minh khả năng hỗ trợ bác sĩ chẩn đoán bệnh nhiễm trùng huyết thông qua
các luật suy diễn logic dạng \code{IF <tiền đề> THEN <kết luận>}.

\begin{dinhnghia}[Hệ chuyên gia dựa trên luật]
Hệ chuyên gia dựa trên luật \en{(Rule-based Expert System)} gồm ba thành phần cơ bản:
\begin{enumerate}[leftmargin=1.5em, itemsep=2pt]
  \item \textbf{Cơ sở tri thức \en{(Knowledge Base)}:} Chứa các sự kiện \en{(facts)}
        và tập luật \en{(rules)} do chuyên gia y tế biên soạn.
  \item \textbf{Động cơ suy diễn \en{(Inference Engine)}:} Áp dụng các thuật toán
        suy diễn xuôi \en{(forward-chaining)} hoặc suy diễn ngược
        \en{(backward-chaining)} để tìm lời giải từ dữ liệu đầu vào.
  \item \textbf{Bộ nhớ làm việc \en{(Working Memory)}:} Lưu trữ trạng thái tạm thời,
        các sự kiện đã được kích hoạt trong phiên làm việc.
\end{enumerate}
\end{dinhnghia}

Ưu điểm tuyệt đối của hệ chuyên gia trong y tế là \textbf{tính tất định
\en{(Determinism)}}, \textbf{tính minh bạch có thể giải thích \en{(Explainability)}}
và \textbf{an toàn tuyệt đối}: một khi tập luật được thẩm định bởi hội đồng y khoa,
hệ thống sẽ không bao giờ tự ý đưa ra kết luận nằm ngoài tri thức đã nạp \cite{Musen2014}.

\section{Phân tầng nguy cơ y tế và giao thức Manchester}

Phân loại bệnh nhân \en{(Medical Triage)} là quy trình sắp xếp thứ tự ưu tiên
chăm sóc y tế dựa trên mức độ nghiêm trọng của bệnh nhân nhằm tối ưu hóa thời
gian cứu chữa và sử dụng nguồn lực hiệu quả.

Hệ thống phân loại Manchester \en{(Manchester Triage System - MTS)}
\cite{Mackway2014} là tiêu chuẩn lâm sàng quốc tế phổ biến hàng đầu. MTS vận
hành theo nguyên tắc tìm kiếm các yếu tố chỉ điểm phân biệt \en{(discriminators)},
đặc biệt là các \textbf{dấu hiệu cờ đỏ \en{(Red Flags)}}. Trong đồ án này, quy
chuẩn phân tầng được chuẩn hóa thành 4 nhóm kết quả:
\begin{itemize}[leftmargin=1.6em, itemsep=2pt]
  \item \term{Cấp cứu \en{(Emergency)}:} Xuất hiện các dấu hiệu đe dọa tính mạng
        như đau ngực lan toả, khó thở cấp, dấu hiệu đột quỵ FAST, sốt cao co
        giật. Bệnh nhân cần gọi cấp cứu 115 ngay lập tức.
  \item \term{Đi khám bác sĩ \en{(Refer Doctor)}:} Tình trạng bệnh dai dẳng hoặc có
        nguy cơ chuyển biến nặng, cần can thiệp cận lâm sàng tại phòng khám.
  \item \term{Tư vấn dược sĩ \en{(Pharmacist Review)}:} Các ca bệnh không rõ ràng,
        bệnh nhân có tiền sử đặc biệt (mang thai, đang dùng nhiều thuốc), cần
        dược sĩ chuyên môn thẩm định trực tiếp.
  \item \term{Tự chăm sóc \en{(Self-care)}:} Các triệu chứng nhẹ, mới khởi phát,
        có thể tự theo dõi và dùng các thuốc không kê đơn \en{(OTC)} thông thường.
\end{itemize}

Bảng~\ref{tab:manchester} tổng hợp bốn mức xử trí cùng thời gian đích tham
chiếu theo tinh thần phân loại Manchester: mức càng nghiêm trọng thì đáp ứng
càng phải tức thì.

\begin{table}[H]
\centering
\caption{Bốn mức phân tầng nguy cơ và thời gian đích tham chiếu}
\label{tab:manchester}
\small
\begin{tabular}{|l|l|p{5.6cm}|}
\toprule
\textbf{Mức xử trí} & \textbf{Thời gian đích} & \textbf{Đặc điểm nhận diện} \\
\midrule
Cấp cứu & Ngay lập tức (115) & Dấu hiệu đe dọa tính mạng: đau ngực dữ dội,
khó thở cấp, đột quỵ, co giật, mất ý thức. \\
Đi khám & Trong ngày & Bệnh dai dẳng, sốt kéo dài, dấu hiệu mất nước nặng,
triệu chứng ở trẻ nhỏ. \\
Tư vấn dược sĩ & Trong ngày tại nhà thuốc & Ca không rõ ràng hoặc có yếu tố
bối cảnh: mang thai, cho con bú, đang dùng nhiều thuốc. \\
Tự chăm sóc & Theo dõi tại nhà & Triệu chứng nhẹ mới khởi phát, đáp ứng tốt
với thuốc OTC và nghỉ ngơi. \\
\bottomrule
\end{tabular}
\end{table}

\section{Các công trình liên quan}

\subsection{Hệ kiểm tra triệu chứng thương mại}

Lĩnh vực sàng lọc triệu chứng tự động đã có nhiều sản phẩm thương mại và nhiều
nghiên cứu đo lường độc lập. Semigran và cộng sự \cite{Semigran2015} kiểm tra
23 hệ kiểm tra triệu chứng trực tuyến trên 45 ca bệnh chuẩn hoá và cho thấy
chỉ 57\% lời tư vấn phân tầng là phù hợp; đáng chú ý, xu hướng chung là
\emph{thiên về bảo thủ} — khuyến cáo người dùng đi khám cả ở những ca vốn có
thể tự chăm sóc. Gilbert và cộng sự \cite{Gilbert2020} so sánh 8 ứng dụng phổ
biến (Ada, Babylon, K Health, Your.MD \dots) với 7 bác sĩ đa khoa trên 200 ca:
không ứng dụng nào vượt được bác sĩ về độ chính xác chẩn đoán, nhưng một số ít
ứng dụng đạt độ an toàn tư vấn ngang bác sĩ. Các tổng quan gần đây cũng nhấn mạnh
rằng độ chính xác và độ an toàn của symptom checker phụ thuộc mạnh vào thiết kế
bộ ca, nhãn chuẩn và cách tính over-triage/under-triage
\cite{RiboliSasco2023, Painter2022}. Bảng~\ref{tab:related} tổng hợp so sánh
với đồ án này.

\begin{table}[H]
\centering
\caption{So sánh các công trình liên quan với hệ thống của đồ án}
\label{tab:related}
\small
\setlength{\tabcolsep}{3.0pt}
\begin{tabular}{|p{2.6cm}|p{3.2cm}|p{3.2cm}|p{4.6cm}|}
\toprule
\textbf{Hệ thống} & \textbf{Kiến trúc} & \textbf{Điểm mạnh} & \textbf{Hạn chế đồ án không kế thừa} \\
\midrule
MYCIN \cite{Buchanan1984} & Hệ chuyên gia luật, suy diễn ngược & Minh bạch,
truy vết được từng luật & Giao diện hội thoại thô; không hiểu ngôn ngữ tự nhiên. \\
Semigran 2015 \cite{Semigran2015} (23 hệ) & Đa dạng: luật hoặc ML & Tận dụng
tri thức y khoa tĩnh & Độ chính xác tư vấn thấp (57\%); không hỗ trợ tiếng Việt. \\
Ada, Babylon \cite{Gilbert2020} & Suy luận xác suất + tri thức lớn & Độ phủ
bệnh học rộng, kiểm chứng lâm sàng & Hộp đen với người dùng; không thể truy
vết về luật cụ thể. \\
LLM y tế trực tiếp \cite{Singhal2023} & Mô hình ngôn ngữ lớn tinh chỉnh & Hiểu
văn phong tự nhiên đa ngôn ngữ & Ảo giác; không xác định được ranh giới an toàn. \\
\midrule
\textbf{Đồ án này} & \textbf{Lai: luật tất định + LLM ngoại vi} & Quyết định
truy vết về luật; kiểm duyệt đầu ra AI; tiếng Việt & Độ phủ tri thức 31 triệu
chứng còn hẹp. \\
\bottomrule
\end{tabular}
\end{table}

\subsection{Điểm khác biệt của đồ án}

Từ đối chiếu trên, đồ án chọn điểm đón đầu khác các hệ đã có: \textbf{(1)}
mọi quyết định sinh tử đều nằm trong động cơ luật minh bạch, mỗi kết luận
truy vết được về \code{rule\_id} cụ thể — khắc phục điểm yếu hộp đen của
nhóm ứng dụng thương mại; \textbf{(2)} mô hình ngôn ngữ lớn chỉ đứng ở hai
vị trí ngoại vi (đọc hiểu đầu vào, diễn giải đầu ra) và bị khoá bằng kiểm
duyệt ở tầng mã nguồn — khắc phục rủi ro ảo giác của hướng LLM thuần;
\textbf{(3)} hướng thẳng người dùng tiếng Việt, vốn là đối tượng chưa có hệ
kiểm tra triệu chứng nào được kiểm chứng bài bản.

\section{Mô hình ngôn ngữ lớn và kỹ thuật Tool-calling}

Các mô hình ngôn ngữ lớn \en{(LLMs)} như GPT-4 hay Claude \cite{Singhal2023}
sở hữu năng lực vượt trội trong việc hiểu văn phong tiếng Việt đời thường,
chuyển đổi ngôn ngữ miêu tả mơ hồ của người bệnh thành các thuật ngữ khoa học.

Tuy nhiên, do bản chất thống kê xác suất tiếp theo của từ \en{(next-token
prediction)}, LLM dễ bị ảo giác \cite{Thirunavukarasu2023}. Để giải quyết hạn
chế này, kỹ thuật \term{Gọi công cụ \en{(Tool/Function Calling)}}
\cite{Schick2023} được áp dụng. Thay vì để LLM tự do sinh văn bản, hệ thống
cung cấp một lược đồ JSON \en{(JSON Schema)} định nghĩa nghiêm ngặt cấu trúc đầu ra:

\begin{vidu}[Lược đồ công cụ trích xuất sự kiện y tế]
Mô hình bị ràng buộc phải trả về các tham số khớp với schema:
\begin{itemize}[leftmargin=1.5em, itemsep=1pt]
  \item \code{symptom\_codes}: Danh sách mã triệu chứng chuẩn hoá (\code{sot},
        \code{dau\_dau}, \dots).
  \item \code{duration\_days}: Số ngày khởi phát (kiểu số nguyên).
  \item \code{has\_red\_flag}: Cờ nhị phân biểu thị sự hiện diện dấu hiệu nguy kịch.
\end{itemize}
\end{vidu}

Nhờ kỹ thuật này, đầu ra của AI trở thành dữ liệu có cấu trúc, cho phép mã nguồn
kiểm định tính hợp lệ \en{(validation)} và đưa vào bộ suy diễn luật một cách tin cậy.

\section{Hệ sinh thái công nghệ và kiến trúc triển khai}

Hệ thống được xây dựng trên nền tảng các công nghệ web hiện đại, ổn định:
\begin{itemize}[leftmargin=1.6em, itemsep=2pt]
  \item \textbf{Laravel 13 Framework (PHP 8.2):} Cung cấp kiến trúc MVC mạnh mẽ,
        hệ thống Object-Relational Mapping (Eloquent ORM) tối ưu, hàng đợi tác
        vụ bất đồng bộ \en{(Database Queue)} và bảo mật cao.
  \item \textbf{MySQL Database:} Quản lý cơ sở dữ liệu quan hệ, hỗ trợ kiểu dữ liệu
        JSON hiện đại và lập chỉ mục tối ưu cho các truy vấn tra cứu luật nhanh.
  \item \textbf{OpenAI-compatible AI Service API:} Giao tiếp linh hoạt với các
        nhà cung cấp mô hình ngôn ngữ lớn thông qua chuẩn API mở, hỗ trợ cơ chế
        chuyển đổi dự phòng \en{(Failover Provider)} khi mất kết nối.
  \item \textbf{Tailwind CSS:} Xây dựng giao diện hướng người dùng theo phong cách
        Tin cậy Y tế \en{(Clinical Trust)}, trực quan và đáp ứng đa thiết bị.
\end{itemize}
